\documentclass[journal]{IEEEtran}

\usepackage[T1]{fontenc}
\usepackage{amsmath,amssymb,amsfonts}
\usepackage{bm}
\usepackage{booktabs}
\usepackage{multirow}
\usepackage{array}
\usepackage{graphicx}
\usepackage{url}
\usepackage{cite}
\usepackage{xcolor}
\usepackage{algorithm}
\usepackage{algpseudocode}

\begin{document}

\title{A Complete Teacher--Student Knowledge Distillation Pipeline for\
Aerial Land Semantic Segmentation on the MBRSC Dubai Dataset}

\author{First~Author,~\IEEEmembership{Member,~IEEE},
        Second~Author,~\IEEEmembership{Member,~IEEE},
        and~Third~Author,~\IEEEmembership{Senior Member,~IEEE}%
\thanks{Manuscript received April 15, 2026.}%
\thanks{First Author and Second Author are with the Department of Computer Vision, XYZ University, City, Country (e-mail: author@xyz.edu).}%
\thanks{Third Author is with the Department of AI Systems, ABC Institute, City, Country.}%
}

\markboth{IEEE Transactions on \LaTeX{} Template,~Vol.~X, No.~X, April~2026}%
{Author \MakeLowercase{\textit{et al.}}: Complete KD Pipeline for Aerial Land Segmentation}

\maketitle

\begin{abstract}
This paper presents an end-to-end semantic segmentation system for aerial land-cover analysis using teacher--student knowledge distillation (KD). The system is implemented as a complete reproducible pipeline: (i) raw-tile preprocessing into class-index masks and training patches, (ii) high-capacity teacher training with a Swin-Base UNet++ architecture, (iii) student distillation with an EfficientNet-B2 UNet architecture, (iv) standardized evaluation with pixel-level metrics, and (v) deployment-oriented inference utilities. The teacher combines hierarchical transformer encoding, Atrous Spatial Pyramid Pooling (ASPP), dense UNet++ decoding, and deep supervision. The student uses a compact encoder--decoder with projection heads for feature-level mimicry. Distillation jointly optimizes hard supervision, soft-logit transfer (temperature-scaled Kullback--Leibler divergence), and multi-level feature regression. The implementation includes class-imbalance correction, mixed-precision training, gradient clipping, learning-rate scheduling, checkpointing, and TensorBoard instrumentation. This manuscript documents all technical details directly from the codebase so the work is publication-ready and fully auditable.
\end{abstract}

\begin{IEEEkeywords}
Semantic segmentation, aerial imagery, knowledge distillation, Swin Transformer, UNet++, EfficientNet, ASPP, deep supervision.
\end{IEEEkeywords}

\section{Introduction}
\IEEEPARstart{A}{erial} semantic segmentation is a core primitive for urban monitoring, infrastructure analytics, water-resource tracking, and environmental planning. In practice, dense prediction quality must be balanced against inference efficiency: large transformer-based models are accurate but expensive, whereas lightweight convolutional models are deployable but less expressive. This work addresses that gap through a complete teacher--student KD pipeline tailored to high-resolution remote-sensing imagery.

The implemented system targets a six-class segmentation problem (building, land, road, vegetation, water, unlabeled) and is organized as a full stack: data preprocessing, dataset handling, model definitions, training scripts, loss engineering, evaluation metrics, visual diagnostics, and inference scripts. The codebase is structured to support both research rigor and production-style reproducibility.

\subsection{Contributions}
The project contributes the following concrete technical elements:
\begin{itemize}
    \item A Swin-Base UNet++ teacher with ASPP bottleneck and deep supervision.
    \item A compact EfficientNet-B2 UNet student with feature projection heads.
    \item A three-term KD objective combining hard labels, soft labels, and feature mimicry.
    \item Class-imbalance handling via inverse-frequency weights computed from training masks.
    \item End-to-end engineering details: augmentation, mixed precision, clipping, scheduling, logging, checkpoint policies, and inference/visualization utilities.
\end{itemize}

\section{Dataset and Problem Setup}
\subsection{Dataset Organization}
The pipeline operates on the MBRSC-style Dubai semantic segmentation dataset, stored as tiled folders (Tile 1 to Tile 8), each with RGB images and RGB masks. The class vocabulary is:
\begin{enumerate}
    \item building
    \item land
    \item road
    \item vegetation
    \item water
    \item unlabeled
\end{enumerate}

\subsection{Color-to-Class Mapping}
Raw masks are RGB-valued and converted to integer labels using nearest-color assignment in Euclidean color space. Let $\mathbf{p}_i\in\mathbb{R}^3$ be pixel color and $\mathbf{c}_k$ class color for class $k\in\{0,\dots,5\}$. The assigned label is
\begin{equation}
\hat{y}_i = \arg\min_k \|\mathbf{p}_i-\mathbf{c}_k\|_2^2.
\end{equation}
This choice is robust to boundary anti-aliasing artifacts compared with strict exact-color matching.

\subsection{Patch Extraction Pipeline}
Each raw image/mask pair is:
\begin{enumerate}
    \item resized to $672\times672$ (divisible by tile size),
    \item transformed to class-index mask,
    \item tiled into non-overlapping $224\times224$ patches,
    \item saved as PNG image patches and NumPy mask arrays.
\end{enumerate}

If tile size is $s=224$ and resized side length is $r=672$, each resized image yields $(r/s)^2=9$ patches.

\subsection{Train/Validation Split}
Patch IDs are split into train/validation sets via random split (default train ratio $0.8$) using a fixed random seed. Split IDs are serialized into text files (one ID per line) and consumed by the dataset loader.

\subsection{Data Normalization and Augmentation}
Image normalization uses ImageNet statistics:
\begin{equation}
\mu=(0.485,0.456,0.406), \quad \sigma=(0.229,0.224,0.225).
\end{equation}
Training augmentations include horizontal flip, vertical flip, random 90-degree rotation, and color jitter. Validation applies only normalization and tensor conversion.

\section{Model Architectures}
\subsection{Teacher: Swin-Base UNet++}
\subsubsection{Encoder}
The teacher encoder is \texttt{swin\_base\_patch4\_window7\_224} (timm, \texttt{features\_only=True}) with outputs from four stages. Channel dimensions are:
\begin{equation}
[128,\;256,\;512,\;1024].
\end{equation}
Feature tensors are converted to $\text{B}\times\text{C}\times\text{H}\times\text{W}$ layout when required.

\subsubsection{ASPP Bottleneck}
The deepest feature map is processed by ASPP with branches: one $1\times1$ conv branch, three $3\times3$ dilated branches (dilations 6, 12, 18), and one global-pooling branch. Concatenated branch outputs are projected to 512 channels.

\subsubsection{UNet++ Dense Decoder}
Decoder nodes use dense skip aggregation across scales. Each node is a residual dense block with mixed dilation (including dilation 2 in intermediate conv) to enlarge receptive field while preserving spatial detail.

\subsubsection{Deep Supervision}
During training, the model returns one main output and three auxiliary outputs from intermediate decoder levels. During evaluation/inference, only the main segmentation logits are returned.

\subsection{Student: EfficientNet-B2 UNet}
\subsubsection{Encoder}
The student encoder is \texttt{efficientnet\_b2} (timm, \texttt{features\_only=True}, \texttt{out\_indices=(1,2,3,4)}). Effective channel dimensions are:
\begin{equation}
[24,\;48,\;120,\;352].
\end{equation}

\subsubsection{Decoder}
A lightweight UNet-style decoder progressively upsamples and fuses skip features. Decoder channel path is approximately
\begin{equation}
(352+120)\to128,\quad(128+48)\to64,\quad(64+24)\to32,\quad32\to32.
\end{equation}
A $1\times1$ head maps to six classes.

\subsubsection{Projection Heads for KD}
To align student features with teacher feature dimensions, $1\times1$ projection heads are used:
\begin{equation}
128\to256,\quad64\to128,\quad32\to64.
\end{equation}
Returned student feature dict keys are \texttt{d3}, \texttt{d2}, \texttt{d1}.

\section{Loss Functions and Distillation Objective}
\subsection{Teacher Segmentation Loss}
Teacher training uses a weighted sum:
\begin{equation}
\mathcal{L}_{\text{teacher-base}}=0.4\mathcal{L}_{\text{CE}}+0.4\mathcal{L}_{\text{Focal}}+0.2\mathcal{L}_{\text{Dice}}.
\end{equation}
If deep supervision outputs are present, each auxiliary output contributes an additional $0.4\,\mathcal{L}_{\text{teacher-base}}$.

\subsection{Cross-Entropy with Class Weights}
Class weights are loaded from \texttt{configs/class\_weights.npy} if available. They are computed from training masks using inverse frequency:
\begin{equation}
w_k = \frac{1}{n_k+1}, \quad
\tilde{w}_k = \frac{6w_k}{\sum_j w_j},
\end{equation}
then clipped to $[0.5,5.0]$.

\subsection{Focal Loss}
For per-pixel CE term $\ell_{\text{CE}}$, with $p_t=\exp(-\ell_{\text{CE}})$ and focusing parameter $\gamma=2$:
\begin{equation}
\mathcal{L}_{\text{Focal}} = (1-p_t)^{\gamma}\,\ell_{\text{CE}}.
\end{equation}

\subsection{Dice Loss}
With class-wise softmax probabilities and one-vs-rest masks:
\begin{equation}
\mathcal{L}_{\text{Dice}}=1-\frac{1}{C}\sum_{c=1}^{C}
\frac{2\sum p_c y_c + \epsilon}{\sum p_c + \sum y_c + \epsilon},
\end{equation}
with $C=6$ and $\epsilon=10^{-6}$.

\subsection{Knowledge Distillation Loss}
Student optimization combines hard, soft, and feature losses:
\begin{align}
\mathcal{L}_{\text{hard}} &= 0.4\mathcal{L}_{\text{CE}} + 0.4\mathcal{L}_{\text{Focal}} + 0.2\mathcal{L}_{\text{Dice}}, \\
\mathcal{L}_{\text{KL}} &= T^2\,\mathrm{KL}\Big(\log\mathrm{softmax}(\tfrac{z_s}{T}),\,\mathrm{softmax}(\tfrac{z_t}{T})\Big), \\
\mathcal{L}_{\text{feat}} &= \frac{1}{|\mathcal{K}|}\sum_{k\in\mathcal{K}}\|f_s^k - \mathrm{Resize}(f_t^k)\|_2^2, \\
\mathcal{L}_{\text{KD}} &= \beta\mathcal{L}_{\text{hard}} + \alpha\mathcal{L}_{\text{KL}} + \gamma\mathcal{L}_{\text{feat}}.
\end{align}
Default distillation hyperparameters are:
\begin{equation}
T=4.0,\quad \alpha=0.5,\quad \beta=0.3,\quad \gamma=0.2.
\end{equation}

\section{Training Pipeline}
\subsection{Teacher Training Protocol}
\textbf{Model:} Swin-Base UNet++ with deep supervision enabled.

\textbf{Optimizer:} AdamW with differential learning rates:
\begin{itemize}
    \item encoder LR = $1\times10^{-5}$
    \item non-encoder LR = $1\times10^{-4}$
    \item weight decay = $1\times10^{-2}$
\end{itemize}

\textbf{Schedule:} linear warmup (10 epochs) followed by cosine decay for total 150 epochs.

\textbf{Batching:} batch size 4 (GPU constrained setting), train/val DataLoaders with \texttt{pin\_memory=True} and persistent workers when workers $>0$.

\textbf{Stability:}
\begin{itemize}
    \item Automatic Mixed Precision (AMP) on CUDA,
    \item gradient clipping with max norm 1.0,
    \item periodic validation every 5 epochs (plus epoch 1).
\end{itemize}

\textbf{Checkpoint policy:}
\begin{itemize}
    \item save best by mIoU to \texttt{teacher\_best.pth},
    \item save best by pixel accuracy to \texttt{teacher\_best\_acc.pth},
    \item save periodic state checkpoints every 25 epochs.
\end{itemize}

\subsection{Student KD Training Protocol}
\textbf{Teacher loading:} teacher checkpoint loaded and frozen (all params \texttt{requires\_grad=False}).

\textbf{Feature hooks:} forward hooks are attached to teacher decoder nodes (\texttt{node\_0\_1}, \texttt{node\_0\_2}, \texttt{node\_0\_3}) to capture intermediate features \texttt{d3}, \texttt{d2}, \texttt{d1}.

\textbf{Student optimization:}
\begin{itemize}
    \item optimizer: AdamW, LR $1\times10^{-3}$, weight decay $1\times10^{-4}$,
    \item scheduler: cosine annealing for 100 epochs,
    \item batch size: 8,
    \item AMP + gradient clipping (max norm 1.0).
\end{itemize}

\textbf{Validation/checkpoint:} validate every 5 epochs (plus epoch 1), save best student by mIoU to \texttt{student\_best.pth}.

\subsection{Implementation-Level Training Algorithm}
\begin{algorithm}[t]
\caption{Student KD Training (per epoch)}
\begin{algorithmic}[1]
\For{each batch $(x,y)$}
    \State Run frozen teacher: $(z_t, f_t)\leftarrow \mathcal{T}(x)$ (features via hooks)
    \State Run student: $(z_s, f_s)\leftarrow \mathcal{S}(x)$
    \State Compute $\mathcal{L}_{\text{KD}}=\beta\mathcal{L}_{\text{hard}}+\alpha\mathcal{L}_{\text{KL}}+\gamma\mathcal{L}_{\text{feat}}$
    \State Backpropagate with AMP scaling
    \State Unscale gradients and apply clip-norm ($1.0$)
    \State Optimizer step and scaler update
\EndFor
\State Step LR scheduler
\State Periodically evaluate mIoU/accuracy on validation loader
\end{algorithmic}
\end{algorithm}

\section{Evaluation Methodology}
\subsection{Primary Metrics}
Given class set $\{0,\dots,C-1\}$:
\begin{align}
\text{IoU}_c &= \frac{TP_c}{TP_c+FP_c+FN_c},\\
\text{mIoU} &= \frac{1}{C}\sum_{c=0}^{C-1}\text{IoU}_c,\\
\text{Pixel Accuracy} &= \frac{\sum_c TP_c}{\text{all pixels}}.
\end{align}

\subsection{Extended Metrics}
The evaluation module also supports:
\begin{itemize}
    \item per-class IoU,
    \item macro F1 score,
    \item streaming confusion matrix accumulation for epoch-level statistics.
\end{itemize}

\subsection{Streaming Confusion Matrix}
Predictions and labels are accumulated into a $C\times C$ matrix without storing full dataset outputs in memory. This supports robust and efficient metric computation on large validation sets.

\section{Comparative Results Against Published Models}
To explicitly compare our implementation with recent literature, we include two paper baselines provided by the user: a hybrid Swin-UNet variant~\cite{sundarr2025hybrid} and an InceptionResNetV2-UNet transfer-learning model~\cite{faska2025inceptionresnet}. For fair reporting, all values should correspond to the same dataset split and identical metric definitions (mIoU, pixel accuracy, and macro F1).

\begin{table}[t]
\caption{Comparison With Published Segmentation Models}
\label{tab:paper_comparison}
\centering
\begin{tabular}{lccc}
\toprule
Model & mIoU (\%) & Accuracy (\%) & Macro F1 (\%) \\
\midrule
Hybrid Swin-UNet (paper)~\cite{sundarr2025hybrid} & -- & -- & -- \\
InceptionResNetV2-UNet (paper)~\cite{faska2025inceptionresnet} & -- & -- & -- \\
\midrule
Our Teacher (Swin-Base UNet++) & -- & -- & -- \\
Our Student (EfficientNet-B2 UNet, KD) & -- & -- & -- \\
\bottomrule
\end{tabular}
\end{table}

In the codebase, comparison plots are now generated with a white background for publication consistency. The plotting script also accepts external paper baseline metrics (JSON input), so the same figure can overlay our teacher/student scores with paper-reported values.

\section{Inference and Visualization Utilities}
\subsection{Model-Agnostic Predictor}
Inference code supports both teacher and student types with fallback loading paths (new architecture first, legacy model second). The predictor:
\begin{enumerate}
    \item converts BGR to RGB,
    \item applies normalization,
    \item runs forward pass,
    \item applies argmax over logits,
    \item computes softmax class probabilities,
    \item produces colored masks and overlay visualizations.
\end{enumerate}

\subsection{Class Palette and Reporting}
A fixed class color palette is used for rendering outputs. Pixel-wise class distribution summaries are computed and printed for each prediction, aiding qualitative sanity checks.

\subsection{Evaluation Visualization Suite}
The visualization module is engineered for publication-style figures including:
\begin{itemize}
    \item per-class IoU bar charts,
    \item confusion matrix heatmaps,
    \item sample prediction galleries,
    \item precision/recall/F1 grouped charts,
    \item class distribution plots,
    \item teacher-vs-student comparison plots.
\end{itemize}

\section{Configuration and Reproducibility Details}
\subsection{Teacher Configuration (from YAML)}
\begin{itemize}
    \item model name: \texttt{SwinUNetPlusPlus}
    \item backbone: \texttt{swin\_base\_patch4\_window7\_224}
    \item classes: 6
    \item pretrained: true
    \item deep supervision: true
    \item epochs: 150
    \item batch size: 4
    \item optimizer: AdamW
    \item encoder LR: $1\times10^{-5}$
    \item decoder LR: $1\times10^{-4}$
    \item weight decay: $1\times10^{-2}$
    \item warmup epochs: 10
    \item grad clip: 1.0
\end{itemize}

\subsection{Student Configuration (from YAML)}
\begin{itemize}
    \item model name: \texttt{EfficientNetUNet}
    \item backbone: \texttt{efficientnet\_b2}
    \item classes: 6
    \item pretrained: true
    \item KD temperature: 4.0
    \item KD weights: $\alpha=0.5$, $\beta=0.3$, $\gamma=0.2$
    \item teacher checkpoint: \texttt{checkpoints/teacher\_best.pth}
    \item epochs: 100
    \item batch size: 8
    \item LR: $1\times10^{-3}$
    \item weight decay: $1\times10^{-4}$
\end{itemize}

\subsection{Dependency Stack}
The implementation relies on PyTorch, torchvision, timm, segmentation-models-pytorch, albumentations, OpenCV, NumPy, Matplotlib, TensorBoard, tqdm, PyYAML, scikit-learn, and Pillow. These are version-pinned with lower bounds in \texttt{requirements.txt}.

\section{Complexity and Deployment Considerations}
The teacher model is high-capacity and suited for best-accuracy operation; the student model is explicitly designed for lower latency and deployment efficiency. Distillation transfers richer supervision to preserve segmentation quality under reduced parameter budget.

In practical deployment, model choice is task-dependent:
\begin{itemize}
    \item teacher for offline/accuracy-critical analysis,
    \item student for real-time or resource-constrained execution.
\end{itemize}

\section{Limitations and Risk Notes}
\begin{itemize}
    \item Domain shift risk: performance may degrade on imagery distributions far from training data.
    \item Label noise sensitivity: color-based mask conversion assumes consistent annotation palette.
    \item Class imbalance remains challenging for rare classes despite weighting and focal modulation.
    \item Hook-based feature extraction requires architectural consistency between teacher definitions and checkpoints.
\end{itemize}

\section{Conclusion}
This work documents a complete, technically explicit IEEE-style pipeline for aerial land semantic segmentation with knowledge distillation. Every major component in the project is specified: preprocessing, dataset handling, architectures, losses, optimization, metrics, inference, and visualization. The resulting framework is both research-grade and implementation-grounded, enabling direct reproduction, extension, and manuscript submission.

\appendices
\section{Codebase Module Map}
\begin{itemize}
    \item \texttt{datasets/preprocess.py}: raw tile processing, color-to-class conversion, tiling, split generation.
    \item \texttt{datasets/mbrsc\_dataset.py}: dataset loader, augmentations, tensor formatting.
    \item \texttt{models/teacher/swin\_unet\_plus.py}: main teacher model.
    \item \texttt{models/teacher/aspp.py}: ASPP bottleneck module.
    \item \texttt{models/student/efficientnet\_unet.py}: student model + projection heads.
    \item \texttt{training/losses.py}: focal, dice, teacher, and KD losses.
    \item \texttt{training/train\_teacher.py}: teacher optimization loop.
    \item \texttt{training/train\_student\_kd.py}: KD training loop with hooks.
    \item \texttt{training/compute\_class\_weights.py}: inverse-frequency weighting.
    \item \texttt{evaluation/metrics.py}: mIoU, accuracy, per-class IoU, F1, confusion matrix.
    \item \texttt{evaluation/visualize.py}: publication-style evaluation plots.
    \item \texttt{inference/predict.py}: single/batch prediction and overlays.
\end{itemize}

\section{Suggested Figure Placeholders for Final Paper}
\begin{itemize}
    \item System overview diagram (data $\rightarrow$ teacher training $\rightarrow$ KD $\rightarrow$ deployment).
    \item Teacher and student architecture schematics.
    \item Qualitative prediction gallery (input/GT/teacher/student).
    \item Per-class IoU comparison chart.
    \item Confusion matrices for teacher and student.
\end{itemize}

\begin{thebibliography}{99}
\bibitem{hinton2015}
G.~Hinton, O.~Vinyals, and J.~Dean, ``Distilling the knowledge in a neural network,'' \emph{arXiv preprint arXiv:1503.02531}, 2015.

\bibitem{swin2021}
Z.~Liu \emph{et al.}, ``Swin transformer: Hierarchical vision transformer using shifted windows,'' in \emph{Proc. ICCV}, 2021.

\bibitem{unetpp2018}
Z.~Zhou, M.~M.~R. Siddiquee, N.~Tajbakhsh, and J.~Liang, ``UNet++: A nested U-Net architecture for medical image segmentation,'' in \emph{Deep Learning in Medical Image Analysis}, 2018.

\bibitem{deeplabv3}
L.-C. Chen, G.~Papandreou, F.~Schroff, and H.~Adam, ``Rethinking atrous convolution for semantic image segmentation,'' \emph{arXiv preprint arXiv:1706.05587}, 2017.

\bibitem{efficientnet}
M.~Tan and Q.~V. Le, ``EfficientNet: Rethinking model scaling for convolutional neural networks,'' in \emph{Proc. ICML}, 2019.

\bibitem{sundarr2025hybrid}
Sundarr \emph{et al.}, ``Enhanced aerial image segmentation via hybrid Swin-UNet with dilated convolutions and multi-scale fusion,'' 2025.

\bibitem{faska2025inceptionresnet}
Z.~Faska \emph{et al.}, ``Satellite imagery semantic segmentation using InceptionResNetV2-UNet transfer learning model,'' in \emph{Proc. 2025 International Conference on Control, Automation and Diagnosis (ICCAD)}, 2025, pp.~1--6, doi: 10.1109/ICCAD64771.2025.11099155.
\end{thebibliography}

\end{document}
